% ============================================================================
%  МЕТОДОЛОГИЧЕСКАЯ СТАТЬЯ. Законность переноса динамических структур.
%  xelatex method.tex  (дважды). Рисунки: fold.png dispersion.png map_fk.png
%  map_diff2d.png hysteresis.png switch.png coexistence.png loc1.png loc2.png
%  loc3.png lam_clean.png  — в одной папке.
% ============================================================================
\documentclass[11pt]{article}
\usepackage{fontspec}
\usepackage{polyglossia}
\setmainlanguage{russian}\setotherlanguage{english}
\setmainfont{Liberation Serif}
\setsansfont{DejaVu Sans}[Scale=0.90]
\setmonofont{DejaVu Sans Mono}[Scale=0.84]
\usepackage{amsmath,amssymb}
\usepackage{geometry}\geometry{a4paper,left=25mm,right=25mm,top=22mm,bottom=22mm}
\usepackage{graphicx,booktabs,array}
\usepackage{xcolor}
\usepackage{fancyhdr}
\usepackage{caption}
\usepackage{titlesec}
\usepackage{enumitem}
\usepackage[hidelinks]{hyperref}

\definecolor{acmag}{HTML}{9E1B45}\definecolor{accyan}{HTML}{1B6E7E}
\definecolor{acink}{HTML}{18202A}\definecolor{acgrey}{HTML}{555F6B}
\color{acink}
\titleformat{\section}{\sffamily\bfseries\large\color{acmag}}{\thesection.}{0.5em}{}
\titleformat{\subsection}{\sffamily\bfseries\normalsize\color{accyan}}{\thesubsection}{0.5em}{}
\titlespacing{\section}{0pt}{15pt}{6pt}\titlespacing{\subsection}{0pt}{10pt}{3pt}
\captionsetup{font={small,it,color=acgrey},labelfont={small,color=acgrey}}
\setlength{\parindent}{5.2mm}\linespread{1.13}
\pagestyle{fancy}\fancyhf{}\renewcommand{\headrulewidth}{0pt}
\fancyfoot[C]{\color{acgrey}\small\thepage}
\fancyfoot[R]{\color{acgrey}\scriptsize Законность переноса}
\newcommand{\ruleline}{{\color{acmag!40}\hrule height 0.7pt}}
\newcommand{\DA}{D_{\!A}}\newcommand{\DB}{D_{\!B}}
% demarcation tags
\newcommand{\tS}{{\small\sffamily\color{accyan}[строгое]}}
\newcommand{\tT}{{\small\sffamily\color{acmag}[структурное]}}
\newcommand{\tI}{{\small\sffamily\color{acgrey}[интерпретативное]}}

\begin{document}
\thispagestyle{fancy}
\begin{center}
{\sffamily\bfseries\fontsize{19}{24}\selectfont\color{acmag}
Когда перенос динамической структуры на социальное законен}\\[2mm]
{\sffamily\large\color{accyan}
Внутренний критерий на примере реакционно-диффузионной модели}\\[5mm]
{\itshape\color{acgrey}Методологическое исследование в жанре «Изоморфизмов»}\\[1mm]
{\small\color{acgrey}\rule{40mm}{0.4pt}}
\end{center}
\vspace{1mm}
\noindent\ruleline
\begin{quotation}\small\itshape\color{acink}
Перенос структур из теории динамических систем на социальные и экологические предметы обыкновенно
обосновывается сходством поведения и потому уязвим для упрёка в физикализме. В работе предлагается
\emph{внутренний} критерий законности такого переноса: вывод, полученный на модели, переносим ровно в той мере,
в какой он устойчив к снятию тех идеализаций, у которых в предметной области нет аналога~— глобального и
мгновенного управления параметрами, адиабатичности, однородности среды, детерминизма начального условия и
специфики механизма. Критерий операционализирует делёзовскую «механизм-независимость» и расширяет её до
независимости от протокола и от однородности. На контролируемой вычислительной серии с моделью Грея–Скотта
критерий проводится через три переноса. Перенос онтологии виртуального (ДеЛанда) проходит проверку, но
оказывается неполным. Перенос необратимости (Шеффер) в исходной форме не проходит и требует переформулировки в
терминах темпа и геометрии фронтов. Перенос примата среды (Беннетт) проходит и усиливается формальным анализом,
но требует дисциплины риторики. Вклад работы отрицательно-дисциплинирующий: критерий не доказывает, что общества
суть динамические системы, а отсекает заведомо незаконные переносы и на одном проработанном примере показывает,
как это делается.
\end{quotation}
\ruleline
\vspace{2mm}

% ============================================================================
\section{Введение: проблема и возможность}
Родство между узорами, которые самопроизвольно возникают в химических, экологических и социальных системах,
давно замечено и давно же вызывает подозрение. С одной стороны, активатор-ингибиторный принцип~— короткодействующая
фасилитация и дальнодействующая конкуренция~— описывает и пятна на шкуре, и сегрегацию городов
\cite{turing1952,krugman1996}, а каноническая последовательность огрубления растительного покрова служит ранним
сигналом экологического коллапса \cite{rietkerk2004,meron2015}. С другой стороны, история социальной мысли
усеяна обломками таких аналогий: органицизм XIX века, «социальная физика», позднейшая волна приложений синергетики
и диссипативных структур к обществу. Всякий раз соблазн один: увидев в социальном поведении форму, знакомую из
физики, объявить физическую структуру объяснением, а часто и нормой. Это и есть физикализм~— не употребление
математики, а подмена: выдача формального сходства за онтологическое тождество и вывод должного из динамического.

Стандартные способы защититься от физикализма внешни: методологи напоминают, что метафора не есть модель, что
корреляция структур не есть механизм, что «поле» Бурдьё \cite{bourdieu1992} не есть поле физика. Эти напоминания
верны, но бессильны как критерий: они не говорят, \emph{какие именно} выводы из данной модели переносимы, а какие
нет. Нужен инструмент, разделяющий внутри одной и той же модели законную часть переноса от незаконной.

Настоящая работа предлагает такой инструмент и показывает его в действии. Ключевая возможность~— в том, что
вычислительная модель, в отличие от природной системы, позволяет ставить контрфактические эксперименты: можно
принудительно снять любую из идеализаций, на которых держится модель, и посмотреть, переживёт ли интересующий нас
вывод это снятие. Идеализации, которых нет в социальной или экологической мишени,~— это как раз те привилегии,
которыми пользуется программист: он меняет параметры глобально, мгновенно, однородно, из точно воспроизводимого
начального состояния и в рамках одного фиксированного механизма. Если вывод устойчив к снятию этих привилегий,
он говорит о структуре пространства возможностей, а не об артефакте вычислительного удобства,~— и тогда его
перенос законен. Если вывод рушится, его нельзя переносить как есть; его надо либо переформулировать в терминах,
которые снятие переживают, либо отбросить.

Мы формулируем этот критерий (раздел~3), предварительно подведя под него минимальный формальный фундамент
(раздел~2), а затем проводим его через три известных философских переноса~— онтологию виртуального ДеЛанды,
критические переходы Шеффера и живую материю Беннетт (раздел~4). В разделе~5 мы предъявляем сильнейшее
встречное возражение~— исторический упрёк в органицизме~— и показываем, чем критерий от него защищает. Всё это
опирается на ранее проведённую вычислительную серию с моделью Грея–Скотта, включая эксперименты по локальности
параметров, детали которых вынесены в сопроводительный технический отчёт.

Заявляемый вклад скромен и потому защитим. Мы не утверждаем, что общества суть реакционно-диффузионные системы.
Мы предлагаем фильтр, отсекающий незаконные переносы, и демонстрируем его на одном полностью просчитанном случае.

% ============================================================================
\section{Родство с философией науки}
Прежде чем формулировать критерий, стоит указать его место среди уже существующих идей, иначе он покажется либо
самоочевидным, либо произвольным. Он не нов в своём корне, но нов в применении, и точная его новизна видна лишь на
фоне трёх линий философии науки.

\subsection{Робастность}
Первая линия~— \emph{анализ робастности}. Левинс \cite{levins1966} сформулировал его афоризмом: если несколько
независимых, по-разному \emph{неверных} моделей сходятся в одном следствии, это следствие, вероятно, отражает не
их общие огрехи, а реальную черту предмета,~— «истина лежит на пересечении независимых лжей». Уимсатт
\cite{wimsatt2007} обобщил это до общей стратегии различения \emph{реального} и \emph{артефактного}: реально то,
что инвариантно относительно множества независимых способов детекции, доказательства или измерения; артефакт
привязан к частному способу. Вайсберг \cite{weisberg2006} различил виды робастности~— параметрическую,
структурную, репрезентационную.

Наш критерий есть \emph{направленная и экспериментальная} форма этого анализа. Направленная~— потому что мы
варьируем не что попало, а именно те идеализации (И1)–(И5), которых нет в мишени переноса: робастность здесь
привязана к конкретному эпистемическому действию (перенос на общество), а не к общей проверке модели.
Экспериментальная~— потому что вычислительная система позволяет варьировать эти идеализации \emph{контрфактически}
и по одной, чего аналитическая робастность Левинса, работающая со сравнением готовых моделей, обычно не даёт. Иными
словами, мы превращаем робастность из свойства, которое констатируют, в процедуру, которую выполняют.

\subsection{Минимальные модели и универсальность}
Вторая линия отвечает на вопрос, \emph{почему} перенос вообще может быть законен, если модель заведомо ложна как
описание общества. Баттерман \cite{batterman2002} показал, что «минимальные модели»~— намеренно нереалистичные~—
объясняют не вопреки своей ложности, а благодаря ей: они схватывают то, что делает целый \emph{класс
универсальности} систем ведущим себя одинаково у критической точки, независимо от микроскопических деталей.
Объяснительно нагружено здесь не описание конкретного механизма, а структура, общая для всех систем класса.

Это ровно тот уровень, на котором наш перенос и претендует быть законным. Активатор-ингибиторная динамика~—
кандидат в такой класс: короткодействующая фасилитация и дальнодействующая конкуренция дают одну и ту же
морфологию у пятен на шкуре, у растительности и (по гипотезе Кругмана) у городов, безотносительно вещества. То,
что переживает снятие механизм-специфичности (И5),~— это и есть содержание уровня универсальности, а не уровня
реализации. Критерий, стало быть, есть способ эмпирически нащупать, \emph{какая часть} вывода принадлежит классу
универсальности, а какая~— частной модели. Заметим сразу и границу: принадлежность к одному морфологическому
классу узоров не гарантирует общего порождающего механизма~— сходство формы может быть и поверхностным
(об этом~— \S6.3).

\subsection{Модели как посредники}
Третья линия~— тезис Морган и Моррисон \cite{morgan1999} о моделях как \emph{автономных посредниках}: модель не
выводится целиком ни из теории, ни из данных, она частично независимый инструмент, и работа с ней~— это
манипулирование инструментом, показания которого затем ещё нужно отнести к миру. Наш критерий адресует именно
последний шаг: когда показания инструмента (карта режимов, гистерезис, пиннинг) говорят о мишени, а когда~— лишь о
самом инструменте и его настройке. Снятие идеализаций~— это способ откалибровать инструмент: отделить в его
показаниях сигнал о структуре возможного от шума собственных удобств.

На фоне этих трёх линий формулируемый ниже критерий~— не изобретение с нуля, а сведение их в одну процедуру,
заострённое под конкретную задачу переноса на социальное и сделанное выполнимым благодаря контролируемому
вычислительному эксперименту.

% ============================================================================
\section{Минимальный формальный аппарат}
Чтобы критерий не повис в воздухе, а разговор о «переносе» имел точный предмет, необходимо иметь под рукой
элементарный анализ той системы, на которой всё проверяется. Мы вводим ровно столько формализма, сколько нужно,
чтобы ключевые утверждения стали \emph{следствиями}, а не наблюдениями, и чтобы стало видно, где кончается
строгое и начинается интерпретация.

\subsection{Уравнения и однородные состояния}
Модель Грея–Скотта \cite{gray1984,pearson1993} на области с диффузией:
\begin{equation}
\partial_t A=\DA\nabla^2 A-AB^2+f(1-A),\qquad
\partial_t B=\DB\nabla^2 B+AB^2-(k+f)B ,
\label{eq:gs}
\end{equation}
где $A$~— субстрат, $B$~— автокатализатор; $f$~— скорость подпитки, $k$~— скорость удаления. Пространственно
однородные стационарные состояния находятся из $-AB^2+f(1-A)=0$ и $AB^2-(k+f)B=0$. Второе даёт либо тривиальное
состояние
\begin{equation}
(A,B)=(1,0),
\label{eq:trivial}
\end{equation}
либо ветвь $AB=(k+f)$, подстановка которой в первое уравнение приводит к квадратному уравнению
$(k+f)B^2-fB+f(k+f)=0$ с действительными корнями лишь при
\begin{equation}
f\ \ge\ 4\,(k+f)^2 .
\label{eq:fold}
\end{equation}
Неравенство \eqref{eq:fold} задаёт \emph{складку существования} нетривиального однородного стационара: ниже неё
(при меньшем $f$) есть только тривиальное состояние \eqref{eq:trivial}. Эта граница нам понадобится немедленно.

\subsection{Линейная устойчивость и условия Тьюринга}
Линеаризуя \eqref{eq:gs} около однородного стационара $(A^*,B^*)$ и подставляя возмущение $\propto e^{\lambda t+iq\cdot x}$,
получаем дисперсионное соотношение
\begin{equation}
\lambda^2-\tau(q^2)\,\lambda+\Delta(q^2)=0,\quad
\tau(q^2)=\operatorname{tr}J-q^2(\DA+\DB),\quad
\Delta(q^2)=\det\!\big(J-q^2D\big),
\label{eq:disp}
\end{equation}
где $J$~— якобиан реакции, $D=\operatorname{diag}(\DA,\DB)$. В нетривиальном стационаре, где $A^*B^*=(k+f)$,
якобиан выписывается явно:
\begin{equation}
J=\begin{pmatrix} -B^{*2}-f & -2(k+f)\\[2pt] B^{*2} & (k+f)\end{pmatrix},\qquad
\operatorname{tr}J=k-B^{*2},\quad \det J=(k+f)\,(B^{*2}-f),
\label{eq:jac}
\end{equation}
так что однородная устойчивость ($\operatorname{tr}J<0,\ \det J>0$) требует $B^{*2}>\max(f,k)$. Диффузионная
неустойчивость возникает, когда $\Delta(q^2)<0$ на некоторой полосе, что при выполнении однородной устойчивости
эквивалентно
\begin{equation}
\DA J_{22}+\DB J_{11}\ >\ 2\sqrt{\DA\DB\,\det J}\ >0,
\qquad\text{т.е.}\quad \DA(k+f)-\DB(B^{*2}+f)>2\sqrt{\DA\DB\det J}.
\label{eq:turingcond}
\end{equation}
Условие \eqref{eq:turingcond} явно показывает роль \emph{отношения} диффузий: левая часть положительна и достаточно
велика лишь при достаточно большом $\DA/\DB$; при его уменьшении неравенство нарушается и узор становится
невозможен. Маргинальный случай (равенство) даёт критическое волновое число и, следовательно, длину волны узора у
порога:
\begin{equation}
q_c^2=\sqrt{\frac{\det J}{\DA\DB}}\ \Longrightarrow\
\lambda_c=\frac{2\pi}{q_c}\ \propto\ (\DA\DB)^{1/4}\,(\det J)^{-1/4}.
\label{eq:scale}
\end{equation}
Из \eqref{eq:disp}–\eqref{eq:scale} следуют два утверждения, к которым мы будем возвращаться. \emph{Во-первых},
масштаб узора задаётся диффузией: при фиксированном $\DA$ имеем $\lambda_c\propto\DB^{1/4}$~— степень\,\emph{\/}$1/4$,
слабее наивного «диффузионного» $\sqrt{\DB}$. \emph{Во-вторых}, само наличие неустойчивости управляется
\emph{отношением} диффузий: маргинальное условие $\DB J_{11}+\DA J_{22}=2\sqrt{\DA\DB\det J}$ при уменьшении
$\DB/\DA$ (активатор диффундирует всё быстрее относительно ингибитора) перестаёт выполняться, и полоса
неустойчивости закрывается (рис.~\ref{fig:disp}). Реакционные параметры $f,k$ входят сюда лишь через $J$ и
стационар: они двигают рабочую точку, тогда как \emph{отношение диффузий отпирает и запирает самую возможность
диффузионной неустойчивости}. Эта асимметрия~— формальное ядро всего дальнейшего.

\begin{figure}[t]\centering
\includegraphics[width=0.66\linewidth]{dispersion.png}
\caption{Дисперсия $\operatorname{Re}\lambda_{\max}(q)$ для активатор-ингибиторного якобиана при трёх отношениях
диффузий. Полоса неустойчивости ($\lambda>0$) открыта при большом отношении и закрывается при малом: среда, а не
реакция, решает, возможен ли узор. Формальная иллюстрация утверждения~\eqref{eq:disp}.}
\label{fig:disp}
\end{figure}

\subsection{Наш режим субкритичен, и это объясняет всё}
Здесь формальный аппарат приносит неожиданную дисциплину. Для реакционных параметров, в которых стоит вся наша
серия (пятна и лабиринты Грея–Скотта), условие \eqref{eq:fold} \emph{не выполнено}: при $f=0.030,\ k=0.057$ имеем
$4(k+f)^2=0.0303>f$, и нетривиального однородного стационара \emph{нет} (рис.~\ref{fig:fold}). Остаётся лишь
тривиальное состояние \eqref{eq:trivial}, якобиан которого $J=\operatorname{diag}(-f,-(k+f))$ отрицательно
определён, так что при всех $q$
\begin{equation}
\lambda_A(q)=-f-\DA q^2<0,\qquad \lambda_B(q)=-(k+f)-\DB q^2<0 .
\end{equation}
Тривиальное состояние линейно устойчиво \emph{при любой} диффузии: классической тьюринговой неустойчивости здесь
нет вовсе. Наши узоры~— не тьюринговы в учебном смысле, а \emph{субкритические}, нелинейные структуры, живущие
как отдельная ветвь и сосуществующие с устойчивым однородным фоном.

\begin{figure}[t]\centering
\includegraphics[width=0.60\linewidth]{fold.png}
\caption{Складка существования \eqref{eq:fold} в плоскости $(f,k)$. Почти все исследованные паттерн-точки лежат
\emph{ниже} складки~— без нетривиального однородного стационара (субкритический режим); лишь точки заливки
попадают в область его существования.}
\label{fig:fold}
\end{figure}

Это единственное формальное наблюдение объясняет разом всю эмпирику нашей серии, ранее описанную порознь.
\emph{Бистабильность и гистерезис}~— прямое следствие субкритичности: паттерн-ветвь сосуществует с всегда
устойчивым тривиальным состоянием, поэтому при одних параметрах живут оба (рис.~\ref{fig:hyst}).
\emph{Провал глобальной нуклеации при успехе локальной}~— тоже: раз однородное состояние линейно устойчиво,
бесконечно малое возмущение затухает, и запустить паттерн способен лишь конечно-амплитудный локальный зародыш,
преодолевший порог. \emph{Трикритическая точка}~— это точка, где ветвь узора рождается от складки уже с конечной
амплитудой (жёстко), а не непрерывно из нуля (мягко). Формальный минимум, таким образом, не украшение, а то, что
связывает наблюдения в один механизм и заодно дисциплинирует расхожий ярлык: говорить здесь о «тьюринговых
паттернах» строго говоря неверно.

\subsection{Демаркация как рабочий инструмент}
Дальше мы будем помечать каждое существенное утверждение одной из трёх меток, и это не украшение, а способ не
соскользнуть в физикализм незаметно для себя. \tS~— выведено из \eqref{eq:gs}–\eqref{eq:scale} или прямо измерено
в контролируемом эксперименте. \tT~— перенесено на другую предметную область по структурной инвариантности, то
есть прошло критерий раздела~3. \tI~— добавлено интерпретацией; такое утверждение может быть осмысленным и даже
плодотворным, но оно \emph{не диктуется} ни математикой, ни критерием и потому не может служить посылкой в
доказательстве. Правило простое: интерпретативное никогда не переходит в строгое по ходу рассуждения.

% (continued in method2.tex)
